\begin{proposition}[Components under trivial torus monodromy]\hypertarget{tag-01G8}{}\label{cdp:st:D.2.1}
\StatementRecord{cdp:st:D.2.1}{proved}{01G8}{\SeriesRef{cdp:st:D.1.2}, \SeriesRef{cdp:st:C.3.4}}

A smooth compact complex threefold with $b_1=b_2=0$ and a holomorphic
connected-fibre map to a curve, with general fibre a two-torus and
trivial rational monodromy, has reducible fibres. More precisely,
for every nonempty set of $s$ removed values containing all critical
values, their reduced fibres have $r=s+3$ irreducible components in
total. Every component of each reducible fibre has no compact
K\"ahler resolution. Consequently no such fibration can have all
its fibre components admitting compact K\"ahler resolutions.
\end{proposition}


\begin{proof}
\begin{enumerate}[label={$\langle1\rangle$\arabic*.},leftmargin=*]
\item 
Pullback on first rational cohomology from the base is injective because
the fibres are connected. Thus the base is $\PP^1$. Trivial monodromy on
the four-dimensional $H_1(F;\Q)$ gives $t'=4$, and the count in
Proposition~\ref{cdp:st:D.1.2} gives $r=s+3$. If all removed fibres were
irreducible one would have $r=s$, which is impossible. The assertion
about each component is Proposition~\ref{cdp:st:C.3.4}, applied to this
threefold. Such a reducible fibre prevents all fibre components from
admitting compact K\"ahler resolutions.

\end{enumerate}
\end{proof}